\subsection{用生成元和关系定义代数}

设 \(\mathbf{F}\) 是 \(\mathbf{A}\) 上的一个代数， \((x_{i})_{i\in I}\) 是 \(\mathbf{F}\) 的元素的一个族。根据第 7 号命题 7，存在唯一的同态 \(f\colon \operatorname{Lib}_{\mathbf{A}}(\mathbf{I})\to \mathbf{F}\) ，使得 \(f(\mathbf{X}_i) = x_i\) 对所有 \(i\in I\) 成立。要使 \(f\) 为满射，必要且充分的条件是 \((x_{i})_{i\in I}\) 是 \(\mathbf{F}\) 的一个生成系。

若 \(\mathbf{U} \in \operatorname{Lib}_{\mathbf{A}}(\mathbf{I})\) ，则 \(f(\mathbf{U})\) 称为 \(\mathbf{F}\) 的、由 \(\mathbf{U}\) 通过把元素 \(x_{i}\) 代入未定元 \(\mathbf{X}_{i}\) 而得到的元素，也称为 \(\mathbf{U}\) 当未定元取值 \(x_{i}\) （就未定元 \(\mathbf{X}_{i}\) 而言）时的值；它通常记作 \(\mathbf{U}((x_{i})_{i \in I})\) ；特别地， \(\mathbf{U}((\mathbf{X}_{i})_{i \in I}) = \mathbf{U}\) 。若 \(\lambda\) 是 \(\mathbf{F}\) 到一个代数 \(\mathbf{F}'\) （ \(\mathbf{A}\) 上的）中的同态，则

\[
\lambda (\mathrm{U} ((x _ {i}) _ {i \in \mathrm{I}})) = \mathrm{U} ((\lambda (x _ {i})) _ {i \in \mathrm{I}}).
\]

特别考虑以下情形： \(\mathrm{F} = \operatorname{Lib}_{\mathbf{A}}(\mathbf{J})\) ，其中 J 是另一个集合；对元素的每一个族 \((\mathrm{H}_{i})_{i \in I}\) （属于 \(\operatorname{Lib}_{\mathbf{A}}(\mathbf{J})\) ）以及每个族 \((y_{j}')_{j \in J}\) （属于 A-代数 \(F'\) ），

\[
\mathrm{U} ((\mathrm{H} _ {i}) _ {i \in I})) ((y _ {j} ^ {\prime}) _ {j \in J}) = \mathrm{U} ((\mathrm{H} _ {i} ((y _ {j} ^ {\prime}) _ {j \in J})) _ {i \in I}).\tag{37}
\]

在上述记号中， \(\operatorname{Lib}_{\mathbf{A}}(\mathbf{I})\) 的每个使得 \(\mathrm{U}((x_{i})_{i\in\mathbf{I}})=0\) （亦即满足 \(f(\mathrm{U})=0\) ）的元素 U，称为族 \((x_{i})_{i\in\mathbf{I}}\) 在 F 中的关系子。由这些元素组成的双边理想 \(\operatorname{Ker}(f)\) 称为 \((x_{i})\) 的关系子理想。

设 \((\mathbf{R}_j)_{j\in J}\) 是 \(\operatorname{Lib}_{\mathbf{A}}(\mathbf{I})\) 的元素的一个族。 \(((x_i)_{i\in I}, (\mathbf{R}_j)_{j\in J})\) 称为代数 \(F\) 的一个表出，如果 \((x_i)_{i\in I}\) 是 \(F\) 的生成系，且 \(\operatorname{Lib}_{\mathbf{A}}(\mathbf{I})\) 的由诸 \(\mathbf{R}_j\) 生成的双边理想等于族 \((x_i)_{i\in I}\) 的关系子理想； \(x_i\) 称为生成元，而 \(\mathbf{R}_j\) 称为该表出的关系子。

现在考虑任意的集合 I 与元素的一个族 \((\mathbf{R}_j)_{j\in J}\) （属于 \(\operatorname{Lib}_{\mathbf{A}}(\mathbf{I})\) ）。 \(\operatorname{Lib}_{\mathbf{A}}(\mathbf{I})\) 关于由族 \((\mathbf{R}_j)\) 生成的双边理想的商代数 E，称为由生成系 I 经关系子族 \((\mathbf{R}_j)_{j\in J}\) 约束而定义的泛代数。显然，若 \(\overline{\mathbf{X}}_i\) 表示 \(\mathbf{X}_i\) 在 E 中的像，

\[
\left(\left(\overline {{\mathbf {X}}} _ {i}\right) _ {i \in I}, \left(\mathbf {R} _ {j}\right) _ {j \in J}\right)
\]

是 E 的一个表出。此外，若 \((x_{i})_{i\in I}\) 是一个代数 F 的元素族，且 \(\mathbf{R}_j((x_i)_{i\in I}) = 0\) 对所有 \(j\in J\) 成立，则存在唯一的同态 \(g:\mathbf{E}\to \mathbf{F}\) ，使得 \(g(\overline{\mathbf{X}}_i) = x_i\) 对所有 \(i\in I\) 成立；要使 \(((x_i)_{i\in I},(\mathbf{R}_j)_{j\in J})\) 是 F 的一个表出，必要且充分的条件是 \(g\) 为双射。

这些评注为下面的语言滥用提供了正当性：不说“ \(((x_i)_{i \in I}, (R_j)_{j \in J})\) 是 F 的一个表出”，也说“F 是由生成元 \(x_i\) 在关系 \(R_j((x_i)_{i \in I}) = 0\) 约束下生成的代数”。当 \(R_j\) 形如 \(P_j - Q_j\) 时，也说“F 是由 \(x_i\) 在关系 \(P_j((x_i)) = Q_j((x_i))\) 约束下生成的代数”。

设 H 是一个集合；我们将说 \(\operatorname{Lib}_{\mathbf{A}}(\mathbf{H})\) 的元素 S 是 A-代数 F 的一个泛关系子，如果 \(\mathrm{S}((x_{h})_{h\in\mathbb{H}})=0\) 对以 H 为指标集的 F 的元素的每个族 \((x_{h})_{h\in\mathbb{H}}\) 成立。

例子。(1) 取 \(\mathbf{H} = \{1,2,3\}\) ；承认

\[
(\mathrm{X} _ {1} \mathrm{X} _ {2}) \mathrm{X} _ {3} - \mathrm{X} _ {1} (\mathrm{X} _ {2} \mathrm{X} _ {3})
\]

作为泛关系子的代数是结合代数。承认 \(X_{1}X_{2}-X_{2}X_{1}\) 作为泛关系子的代数是交换代数。承认泛关系子 \(X_{1}X_{1}\) 与

\[
(\mathrm{X} _ {1} \mathrm{X} _ {2}) \mathrm{X} _ {3} + (\mathrm{X} _ {2} \mathrm{X} _ {3}) \mathrm{X} _ {1} + (\mathrm{X} _ {3} \mathrm{X} _ {1}) \mathrm{X} _ {2}
\]

作为泛关系子的代数是李代数。

设 I 是一个集合；给定元素的一个族 \((\mathrm{S}_{k})_{k\in\mathbb{K}}\) （属于 \(\operatorname{Lib}_{\mathbf{A}}(\mathbf{H})\) ），考虑 \(\operatorname{Lib}_{\mathbf{A}}(\mathbf{I})\) 中形如 \(\mathrm{S}_{k}((\mathrm{U}_{h}))_{h\in\mathbb{H}}\) 的元素组成的集合 T，其中 k 遍历 K，且对每个 \(k,\left(\mathrm{U}_{h}\right)_{h\in\mathbb{H}}\) 遍历 \(\operatorname{Lib}_{\mathbf{A}}(\mathbf{I})\) 的元素的以 H 为指标集的族的集合；再考虑一个以 T 为元素集的族 \((\mathrm{R}_{j})_{j\in\mathbb{J}}\) 。然后令 F 为由生成系 I 经关系子族 \((\mathrm{R}_{j})_{j\in\mathbb{J}}\) 约束而定义的泛代数，并设 \(u:\operatorname{Lib}_{\mathbf{A}}(\mathbf{I})\to\mathbf{F}\) 为典范同态，于是 \(\operatorname{Ker}(u)\) 由形如 \(\mathrm{S}_{k}((\mathrm{U}_{h})_{h\in\mathbb{H}})\) 的元素生成，其中 \(k\in K\) 且 \((\mathrm{U}_{h})_{h\in\mathbb{H}}\) 是 \(\operatorname{Lib}_{\mathbf{A}}(\mathbf{I})\) 的元素的族；显然每个 \(S_{k}\) \((k\in\mathbb{K})\) 都是 F 的一个泛关系子。现在设 \(F'\) 是承认一个生成系 \((x_{i})_{i\in\mathbb{I}}\) 的代数，对它每个 \(S_{k}\) 都是泛关系子，并设 \(u':\operatorname{Lib}_{\mathbf{A}}(\mathbf{I})\to\mathbf{F}'\) 是使得 \(u'(\mathrm{X}_{i})=x_{i}\) （对所有 \(i\in I\) ）成立的同态；显然 \(\operatorname{Ker}(u)\subset\operatorname{Ker}(u')\) ，因此 \(u'\) 可唯一地写成 \(u'=h\circ u\) 的形式，其中 \(h:F\to F'\) 是一个使 \(h(\overline{\mathrm{X}}_{i})=x_{i}\) 对所有 \(i\in I\) 成立的同态。因此，F 称为由生成系 I 定义的、对应于泛关系子族 \((\mathrm{S}_{k})_{k\in\mathbb{K}}\) 的泛代数。作为语言上的滥用，F 有时也称为由 I 生成且满足恒等式 \(\mathrm{S}_{k}((u_{h}))=0\) （对 F 的元素的每个族 \((u_{h})_{h\in\mathbb{H}}\) ）的泛代数。

例 2. 设 \(\mathbf{L}'\) 为由 \(\mathbf{I}\) 生成且满足恒等式 \((uv)w - u(vw) = 0\) （对 \(\mathbf{L}'\) 的元素的每个三元族）的泛代数，并设 \(\mathbf{L}''\) 为由 \(\mathbf{L}'\) 添加单位元而得到的代数；则存在唯一的保单位同构 \(g\) ，它由 \(\mathbf{L}''\) 到 \(\operatorname{Libas}_{\mathbf{A}}(\mathbf{I})\) 上，使得 \(g(\overline{\mathbf{X}}_i) = \mathbf{X}_i\) 对所有 \(i \in \mathbf{I}\) 成立。因为显然 \(\mathbf{L}''\) 是结合的，且同态 \(g\) 的存在性由 \(\mathbf{L}'\) 的定义及其前面的评注得出；但于是显然 \(\mathbf{L}''\) 满足刻画 \(\operatorname{Libas}_{\mathbf{A}}(\mathbf{I})\)

的泛性质（第 7 号，命题 7），由此得出结论。与上述类似的考虑可应用于结合代数（相应地，交换结合代数），同时考虑到以下评注。当上下文充分表明所考虑的代数是含单位元代数时，作为语言上的滥用，一个族 \((x_{i})_{i\in I}\) （它是代数 \(F\) 的元素的族），若由诸 \(x_{i}\) ( \(i\in I\) ) 与单位元生成的子代数等于 \(F\) ，则常称为 \(F\) 的一个生成系。然后设 \(F\) 是 \(A\) 上的一个含单位元结合代数， \((x_{i})_{i\in I}\) 是 \(F\) 的元素的一个族；根据第 7 号命题 7，存在唯一的保单位同态 \(f:\mathrm{Libas}_{\mathbf{A}}(\mathbf{I})\to F\) ，使得 \(f(\mathbf{X}_i) = x_i\) 对所有 \(i\in I\) 成立；若 \(\mathbf{U}\in \mathrm{Libas}_{\mathbf{A}}(\mathbf{I}), f(\mathbf{U})\) 也称为 \(F\) 的、由 \(\mathbf{U}\) 通过把元素 \(x_{i}\) 代入未定元 \(\mathbf{X}_i\) 而得到的元素，并且也记作 \(\mathrm{U}((x_{i})_{i\in I})\) 。于是，关系子、表出与泛关系子的概念可立即搬到结合代数上；只需把 \(\mathrm{Lib}_{\mathbf{A}}(\mathbf{I})\) 处处替换为 \(\mathrm{Libas}_{\mathbf{A}}(\mathbf{I})\) 即可。由生成系 \(I\) 经关系子族 \((\mathbb{R}_j)_{j\in J}\) 约束而定义的泛含单位元结合代数是 \(\mathrm{Libas}_{\mathbf{A}}(\mathbf{I})\) 关于由族 \((\mathbb{R}_j)\) 生成的双边理想的商代数。由生成系 \(I\) 生成的、对应于一族泛关系子的泛含单位元结合代数类似地定义。我们把关于交换结合代数的类似定义的陈述留给读者，只需以 \(\mathrm{Libasc}_{\mathbf{A}}(\mathbf{I})\) 替换 \(\mathrm{Libas}_{\mathbf{A}}(\mathbf{I})\) 即可。

例 3. 设 L' 是由 I 生成且满足恒等式 uv - vu = 0（对 L' 的元素的每个二元族）的泛含单位元结合代数。如例 2 中那样可以看出，L' 典范同构于 \(\text{Libasc}_{\mathbf{A}}(\mathbf{I})\) 。

